\cleardoublepage

\section{资源受限 MoE 推理中的专家缓存问题建模}

本章定义本文研究的资源受限 MoE 专家缓存问题。第 2 章已经说明，MoE 模型的专家权重访问由路由器动态决定；当全部专家无法常驻 GPU 时，推理过程中的专家访问会转化为缓存命中、专家换入和跨设备传输问题。为了在后续章节中清晰描述 WREC 策略和实验结果，本章先给出问题场景、路由事件表示、专家引用序列、缓存状态、传输代价和瓶颈筛查方法。

需要说明的是，本文并不把硬件瓶颈分析作为新的性能建模框架提出。MoE-Lightning 已经提出 HRM（Hierarchical Roofline Model），用于从 CPU/GPU 计算能力、内存带宽、CPU-GPU 传输带宽和策略变量等方面分析内存受限 MoE 推理的瓶颈~\cite{cao2025moelightning}。本文参考该类分层 Roofline 建模思想，将其用于筛选后续实验所关注的专家传输受限场景（专家传输受限 scenario）；本文的重点是，在这类场景中基于真实路由访问轨迹定义并评估专家缓存策略。

\subsection{问题场景与研究对象}

本文考虑单机资源受限 MoE 推理场景。系统包含 GPU 显存和 CPU 内存两级主要存储资源，模型的非专家参数、KV cache、中间激活和部分专家权重占用 GPU 显存，其余专家权重存放在 CPU 内存中。当某次 MoE 层计算需要访问的专家已经驻留在 GPU 上时，系统可以直接执行专家前馈网络（Feed-Forward Network, FFN）；当该专家不在 GPU 上时，系统需要将专家权重从 CPU 内存加载到 GPU 显存，并在缓存容量不足时换出已有专家。

本文研究对象是给定模型路由结果后的专家权重驻留问题，而不是修改模型结构或改变路由器输出。也就是说，输入词元被路由到哪些专家由原模型决定，本文只讨论在有限 GPU 专家缓存容量下，哪些专家应当保留、哪些专家需要换入，以及缓存策略如何影响未命中次数和专家传输开销。该设定使本文能够在不改变模型输出分布的前提下，隔离分析系统层专家缓存策略的贡献。

\begin{figure}[htbp]
    \centering
    \includegraphics[width=0.51\linewidth]{gpu_vs_cpu_memory_for_moe_inference.png}
    \caption{资源受限 MoE 推理中的专家缓存与权重卸载场景}
    \label{fig:moe-expert-cache-offloading}
\end{figure}

在推理阶段上，本文主要围绕预填充阶段的路由轨迹展开分析。预填充阶段需要对输入提示词中的多个词元进行并行计算，能够形成较长的词元—层级专家访问序列，因此适合用于分析专家访问局部性并开展缓存重放实验。近期 MoE serving 研究也指出，预填充阶段和解码阶段在专家激活密度、请求并发结构和缓存调度目标上存在差异，因而两阶段采用完全相同的专家调度策略往往并非最优~\cite{zhang2025duoservemoe}。解码阶段逐词元生成，访问序列受到自回归依赖和请求调度影响，本文不对连续批处理场景下解码阶段的专家缓存策略作完整建模，而将其作为后续扩展方向讨论。因此，本文结论主要针对预填充阶段的专家访问轨迹。

\subsection{路由事件表示}

MoE 层的路由器为每个词元 选择 top-$k$ 个专家。为了将模型运行过程转化为可分析的系统事件，本文将一次专家选择表示为路由事件：

\begin{equation}
    \label{equ:router-event}
    e_t=(q_t,p_t,l_t,j_t,x_t),
\end{equation}

其中，$q_t$ 表示请求编号，$p_t$ 表示该请求中的 词元位置，$l_t$ 表示 MoE 层编号，$j_t$ 表示 top-$k$ 路由结果中的排序位置，$x_t$ 表示被选中的专家编号。若模型在每个 MoE 层采用 top-$2$ 路由，则同一词元 在同一层会产生两个路由事件。

路由事件保留了本文后续策略需要使用的上下文信息。请求编号和 词元位置用于分析请求内局部性；层编号用于分析跨层专家转移；专家编号用于统计全局热度和缓存命中；top-$k$ 排序位置用于区分同一词元 在同一层的多个专家选择。相比只记录专家编号，路由事件能够更完整地描述 MoE 访问序列的结构。

\subsection{专家引用序列}

为了进行缓存重放，本文将按时间顺序排列的路由事件进一步转化为专家引用序列：

\begin{equation}
    \label{equ:expert-reference-sequence}
    A=(a_1,a_2,\ldots,a_T),
\end{equation}

其中，$a_t=(l_t,x_t)$ 表示第 $t$ 次专家访问对应的层编号和专家编号，$T$ 表示访问总次数。将层编号纳入引用对象，是因为不同 MoE 层中的同名专家并不共享权重。例如，第 3 层的第 1 个专家和第 4 层的第 1 个专家应视为两个不同缓存对象。因此，本文中的“专家对象”指某一 MoE 层中的某一专家，而不是全模型范围内的专家编号。

在该表示下，真实推理过程中的专家权重访问被抽象为一串缓存对象访问。后续重放模拟器在同一条引用序列上运行不同策略，并统计各策略的命中率、传输量和传输时间代理指标。这样可以避免因请求顺序、模型输出或运行时调度差异导致的比较偏差。

\begin{figure}[htbp]
    \centering
    \includegraphics[width=0.51\linewidth]{route_event_to_expert_reference_process.png}
    \caption{路由事件到专家引用序列的转换过程}
    \label{fig:route-event-to-expert-reference}
\end{figure}

\subsection{专家缓存模型}

设 GPU 中可用于专家权重驻留的缓存容量为 $C$。在专家大小近似相同的情况下，$C$ 可以用可容纳的专家对象数量表示；若专家大小不同，则可将容量扩展为字节级约束。本文实验对象中同一模型内各专家 FFN 结构相同，因此后续主要采用专家对象数量作为容量单位。实际系统中的显存碎片、框架开销和对齐开销通过显存余量项近似处理。

记第 $t$ 次访问前的 GPU 专家缓存状态为 $S_t$。当 $a_t\in S_t$ 时，该访问命中缓存；当 $a_t\notin S_t$ 时，该访问发生未命中，系统需要加载专家 $a_t$。若此时缓存已满，则缓存策略需要从 $S_t$ 中选择一个或多个专家换出。缓存状态更新可抽象表示为：

\begin{equation}
    \label{equ:cache-state}
    S_{t+1}=U(S_t,a_t,\pi),
\end{equation}

其中，$U(\cdot)$ 表示缓存状态更新过程，$\pi$ 表示具体缓存策略。当 $a_t\in S_t$ 时，$U$ 仅更新策略 $\pi$ 所维护的访问信息；当 $a_t\notin S_t$ 且 $|S_t|<C$ 时，$U$ 将 $a_t$ 加入缓存；当 $a_t\notin S_t$ 且 $|S_t|=C$ 时，策略 $\pi$ 需要选择待换出的专家对象，然后将 $a_t$ 加入缓存。不同策略的差异主要体现在未命中时的准入、换出和预取决策上。LRU 根据最近访问时间选择换出对象，静态热度策略根据历史访问频率选择驻留对象，WREC 则在第 4 章中进一步结合工作负载统计、请求内局部性和跨层转移信息构造在线评分。

\subsection{代价指标}

本文使用缓存未命中数、未命中率、专家加载次数、专家传输量和传输时间代理指标衡量专家缓存策略。设 $I_t$ 表示第 $t$ 次访问是否未命中：

\begin{equation}
    \label{equ:miss-indicator}
    I_t=
    \begin{cases}
        1, & a_t\notin S_t,\\
        0, & a_t\in S_t.
    \end{cases}
\end{equation}

则总未命中数和未命中率分别为：

\begin{equation}
    \label{equ:miss-rate}
    N_{\mathrm{miss}}=\sum_{t=1}^{T} I_t,\quad
    R_{\mathrm{miss}}=\frac{N_{\mathrm{miss}}}{T}.
\end{equation}

若每个专家对象大小为 $B_{\mathrm{expert}}$，且每次加载均需要从 CPU 向 GPU 传输完整专家权重，则专家传输量可表示为：

\begin{equation}
    \label{equ:transfer-bytes}
    B_{\mathrm{transfer}}=N_{\mathrm{load}}\cdot B_{\mathrm{expert}},
\end{equation}

其中，$N_{\mathrm{load}}$ 表示专家总加载次数。在无预取策略中，专家加载通常由缓存未命中触发，因此 $N_{\mathrm{load}}=N_{\mathrm{miss}}$。在包含预取的策略中，部分专家可能在实际访问前被加载，从而降低后续未命中数；同时，错误预取也会增加额外传输量。因此，需要同时统计未命中数、总加载次数和预取浪费。若记需求加载次数为 $N_{\mathrm{demand}}$，预取加载次数为 $N_{\mathrm{prefetch}}$，则有：

\begin{equation}
    \label{equ:load-count}
    N_{\mathrm{load}}=N_{\mathrm{demand}}+N_{\mathrm{prefetch}}.
\end{equation}

预取浪费可进一步定义为被预取但在换出前未被实际访问的专家对象数量。该指标用于衡量预取策略是否引入了不必要的数据搬移和缓存污染。

为了在基于访问轨迹的重放实验中近似反映专家传输对推理开销的影响，本文使用传输时间代理指标。设 CPU-GPU 有效带宽为 $b_{\mathrm{cg}}$，单位为 Byte/s，则由专家加载引起的传输时间代理可写为：

\begin{equation}
    \label{equ:transfer-time-proxy}
    T_{\mathrm{transfer}}=\frac{B_{\mathrm{transfer}}}{b_{\mathrm{cg}}}.
\end{equation}

该指标表示在不考虑计算与传输重叠时由专家加载引入的传输时间估计，而不直接等同于真实运行时停顿时间。真实系统中还存在计算与传输重叠、调度开销和内存分配开销。因此，在固定模型、固定访问轨迹和固定硬件带宽假设下，本文将该指标作为比较不同专家缓存策略的统一代理指标。

\subsection{基于 HRM 思想的瓶颈筛查}

专家缓存优化并非在所有配置下都有意义。如果全部专家能够常驻 GPU，则专家缓存未命中不会成为主要问题；如果专家传输可以完全隐藏在计算之后，减少未命中对端到端性能的影响也可能有限。因此，在设计和评估缓存策略前，需要先判断目标配置是否处于专家传输受限状态。

本文参考 MoE-Lightning HRM 的分层 Roofline 分析思路，将硬件、模型、工作负载和放置策略共同纳入瓶颈筛查。完整 HRM 会建模 GPU/CPU 计算屋顶、GPU/CPU 内存带宽、CPU-GPU 传输带宽以及 批次大小、micro-批次大小、权重放置比例和 KV cache 放置比例等策略变量~\cite{cao2025moelightning}。本文的筛查目标更窄：判断后续 WREC 实验所采用的模型和硬件配置中，专家权重传输是否可能成为主要瓶颈。

为此，本文首先估算可用于专家缓存的 GPU 显存预算：

\begin{equation}
    \label{equ:expert-cache-memory}
    M_{\mathrm{cache}}=M_{\mathrm{gpu}}-M_{\mathrm{base}}-M_{\mathrm{kv}}-M_{\mathrm{reserve}},
\end{equation}

其中，$M_{\mathrm{gpu}}$ 表示可用 GPU 显存，$M_{\mathrm{base}}$ 表示非专家模型权重及固定运行时状态占用，$M_{\mathrm{kv}}$ 表示 KV cache 预留，$M_{\mathrm{reserve}}$ 表示临时激活、显存碎片和未显式建模开销的安全余量。上述各项在估算时应避免重复计入。由此可以得到专家缓存容量：

\begin{equation}
    \label{equ:expert-cache-capacity}
    C=\left\lfloor\frac{M_{\mathrm{cache}}}{B_{\mathrm{expert}}}\right\rfloor.
\end{equation}

当 $C$ 小于模型中需要频繁访问的专家对象数量时，系统必须在推理过程中进行专家换入和换出。进一步结合 CPU-GPU 带宽 $b_{\mathrm{cg}}$ 和单个专家大小 $B_{\mathrm{expert}}$，可以估算一次专家加载的传输时间：

\begin{equation}
    \label{equ:expert-load-time}
    T_{\mathrm{load}}=\frac{B_{\mathrm{expert}}}{b_{\mathrm{cg}}}.
\end{equation}

若在目标批次大小、提示词长度和缓存容量下，专家加载时间相对 MoE 层计算时间不可忽略，且策略枚举或重放结果显示未命中会带来大量 CPU-GPU 传输，则该配置可视为本文关注的专家传输受限场景。相反，如果全专家驻留可行，或专家传输量相对计算时间很小，则该配置不适合作为验证专家缓存策略收益的主要实验对象。

这种筛查方法的作用是限定本文结论适用范围。本文并不声称所有 MoE 模型都同样适合专家缓存优化，而是在已有 HRM 思想指导下，选择专家传输可能成为瓶颈的配置，再进一步研究真实路由轨迹中是否存在可被在线策略利用的访问局部性。

\subsection{问题定义}

基于以上定义，本文研究的问题可以表述如下：给定 MoE 模型产生的专家引用序列 $A$、专家缓存容量 $C$、专家对象大小 $B_{\mathrm{expert}}$ 和 CPU-GPU 有效带宽 $b_{\mathrm{cg}}$，设计一个在线缓存策略 $\pi$，使其在不使用未来真实访问的条件下，根据当前和历史路由事件维护 GPU 专家缓存状态 $S_t$。该策略的目标是在给定缓存容量约束下，降低缓存未命中数、专家传输量和传输时间代理指标。

本文不直接求解全局最优策略，而是设计可在线执行的启发式策略，并通过与 LRU、静态热度策略和 Belady 离线上界比较评估其效果。该问题与经典缓存问题相似，但具有 MoE 推理场景中的特殊结构。第一，缓存对象由层编号和专家编号共同决定，而不是单一对象编号。第二，访问序列由模型路由器产生，包含请求内 词元顺序、层级结构和 top-$k$ 选择信息。第三，缓存收益受到硬件资源约束影响，只有在专家传输可能成为瓶颈的配置中，减少未命中才更可能转化为实际性能收益。

在后续章节中，本文将策略分为三类进行比较。第一类是在线基线，例如 LRU，只使用当前和历史访问信息。第二类是静态先验策略，例如基于历史工作负载的专家热度常驻。第三类是离线上界，例如 Belady 离线上界策略。该策略在给定引用序列和等大小缓存对象假设下使用未来访问信息选择换出对象，用于衡量在线替换策略与离线最优替换结果之间的差距~\cite{belady1966study}。WREC 属于在线策略，其目标是在不使用未来真实访问的前提下，比简单在线基线更充分地利用 MoE 路由轨迹的结构性。

上述建模将本文的研究对象限定为专家传输受限场景中的在线专家缓存决策：模型路由结果保持不变，策略只根据历史与当前可观测路由事件维护专家驻留集合。基于该问题定义，第 4 章将进一步介绍 WREC 的策略设计、局部性统计方法和重放评估流程。
